\chapter*{Abstract}
\noindent
Automated intrusion response is rarely switched on in operational technology, and the reason is trust rather than detection. A reactive defence that blocks traffic can cut the very process it is meant to protect, and an operator will not arm a measure that might trip a running plant. This project addresses that barrier with CARS, a criticality- and operation-aware software-defined-networking intrusion-response system for industrial control systems. Instead of acting uniformly, CARS conditions its response on two properties that plain blocking ignores: the criticality of the asset under attack, and the industrial operation, recovered from the traffic by a deep-packet-inspection sensor, so that reading a value from a controller is treated differently from writing one. Its active enforcement responses are graded across a seven-rung ladder, each bounded in time and reversible, while the safety-critical loop is held on an alert-only rung that never cuts the process; every decision is recorded.

CARS was built on a physical hardware-in-the-loop testbed rather than a pure network simulation. Two Siemens S7-1212C programmable logic controllers, their operator panels and the software-defined fabric are real hardware running live S7comm and Modbus/TCP, and one of them drives a tank-level process that is co-simulated in Factory~IO to avoid a laboratory water hazard, with an engineering workstation present as in a small plant. The system is realised as a three-table OpenFlow pipeline: identity binding against spoofing, a stateful policy stage carrying the criticality- and operation-aware rulebook, and a learning switch, supported by a flow-integrity self-check and an authenticated control interface. Detection is delegated to the packet sensor, while the software-defined network layer makes the graded decision and enforces it as OpenFlow rules. Every capability was validated at the level of individual packets, protected against unprotected.

The central result is that the armed defence refused every illegitimate operation while cutting none of the legitimate process traffic: a false-positive rate of zero on the live process, no illegitimate operation permitted across the decision space and a 2{,}078-case corpus, with enforcement shown on the wire rather than inferred from a log, and a reaction window with a median of 7.6 milliseconds. A false-data injection that overflowed the tank when unprotected could not move the process when the defence was armed: its S7 session was severed at the first write, at most a single frame reaching the PLC before the cut, and the tank held its band. The evaluation carries its boundaries honestly, including the trusted-insider case that a network layer must not act on. The finding is that the trust barrier which keeps such defences switched off is not fundamental: a reactive network defence can be made criticality- and operation-aware, bounded and reversible enough for an operator to arm on a running process.
